\documentclass[10pt,a4paper]{amsart}
\usepackage{amsmath,amssymb,amsthm,graphicx}
\usepackage{latexsym}
\usepackage{verbatim}
\usepackage{amsmath}
\usepackage{amssymb}
\usepackage{mathrsfs}
\usepackage{amsthm}
\usepackage{tikz}
%\usetikzlibrary{matrix}


\usetikzlibrary{decorations.pathreplacing}
\usetikzlibrary{arrows.meta}





\newcommand{\Ind}{
 \setbox0=\hbox{$x$}\kern\wd0\hbox to 0pt{\hss$
 \mid$\hss}\lower.9\ht0\hbox to 0pt{\hss$\smile$\hss}\kern\wd0
}
\newcommand{\indep}[3]{
 #1\mathop{\mathpalette\Ind{}}_{#2}#3
}
\newcommand{\Indep}{\indep{}{}{}}


\newcommand{\Notind}{
 \setbox0=\hbox{$x$}\kern\wd0\hbox to 0pt{\mathchardef
 \nn=12854\hss$\nn$\kern1.4\wd0\hss}\hbox to 0pt{\hss$\mid$\hss}\lower.9\ht0
 \hbox to 0pt{\hss$\smile$\hss}\kern\wd0
}
\newcommand{\depen}[3]{
 #1\mathop{\mathpalette\Notind{}}_{#2}#3
}



\usepackage{amsthm, amsfonts}
\usepackage{euscript}
\usepackage{graphicx}
\usepackage{float}
\usepackage{caption}
\usepackage{tikz-qtree}
%\usetikzlibrary{arrows.meta}
\usepackage{subcaption}
%%To number theorems better in the introduction
\newtheorem{thrm}{Theorem}
%\renewcommand\thethrm{\Alph{thrm}}

\newcommand{\g}{\gamma}
\newcommand{\G}{\Gamma}
\renewcommand{\d}{\delta}
\renewcommand{\a}{\alpha}
\renewcommand{\epsilon}{\varepsilon}
\renewcommand{\b}{\beta}
\renewcommand{\H}{\mathbb{H}}
%\renewcommand{\dh}{$\delta$-hyperbolic}
\newcommand{\inv}{^{-1}}
\newcommand{\C}{\mathcal{C}}
\newcommand{\K}{\mathcal{K}}
\newcommand{\A}{\mathcal{A}}
\newcommand{\E}{\mathcal{E}}
\newcommand{\cl}{\mathrm{cl}}
\newcommand{\Q}{\mathbb{Q}}
\newcommand{\conv}{\mathrm{conv}}
\newcommand{\acl}{\mathrm{acl}}
\newcommand{\dcl}{\mathrm{dcl}}

\newcommand{\val}{\mathrm{val}}
\newcommand{\tp}{\mathrm{tp}}


\newtheorem{defi}{Definition}[section]
\newtheorem{theorem}[defi]{Theorem}
\newtheorem{definition}[defi]{Definition}
\newtheorem{lemma}[defi]{Lemma}
\newtheorem{proposition}[defi]{Proposition}
\newtheorem{prop}[defi]{Proposition}
\newtheorem{conjecture}[defi]{Conjecture}
\newtheorem{corollary}[defi]{Corollary}
\newtheorem{problem}[defi]{Problem}
\newtheorem{remark}[defi]{Remark}
\newtheorem{example}[defi]{Example}
\newtheorem{examples}[defi]{Examples}
\newtheorem{question}[defi]{Question}
\newtheorem{convention}[defi]{Conventions}
\newtheorem{fact}[defi]{Fact}
\def\Aut{\mathop{\rm Aut}\nolimits}
\def\Isom{\mathop{\rm Isom}\nolimits}
\def\Min{\mathop{\rm Min}\nolimits}
\def\Max{\mathop{\rm Max}\nolimits}
\def\Sym{\mathop{\rm Sym}\nolimits}
\def\nil2{\mathop{\rm nil2}\nolimits}
\def\Th{\mathop{\rm Th}\nolimits}
\def\exp{\mathop{\rm exp}\nolimits}
\def\sup{\mathop{\rm sup}\nolimits}
\def\h{\mathop{\rm ht}\nolimits}
\def\Hom{\mathop{\rm Hom}\nolimits}

\title{Models of the Farey graph theory are trees of Farey pieces}
\author{Zahra Mohammadi Khangheshlaghi, Katrin Tent}
\date{arxiv:2503.02121v2; CC BY 4.0}
\begin{document}
\maketitle
\noindent\textit{Editorial note.} Original complete argument retained with version/license/figure dependencies; all support belongs to this single counted target.\par\medskip
\noindent\textit{Editorial note.} Public Morley rank and current quantifier elimination/forking tasks have different targets. This structural classification is proved by the earlier graph decomposition and reconstruction, which are included rather than assumed.\par\medskip
\noindent\textit{Allowed standard prerequisites (editorial).}\begin{itemize}
\item Standard subject prerequisites. Standard named subject results cited by the authors may be used; all paper-specific supporting arguments are retained verbatim. No modified assumptions or replacement proof.
\end{itemize}\medskip
\section{Introduction}



Our purpose in this article is to investigate the model-theoretic properties of the Farey graph, 
a combinatorial and geometric structure that arises in a number of mathematical fields. 
In geometric group theory, the Farey graph appears e.g. in the study of the free factor complex and the curve graph.

The Farey graph is isomorphic to the complex of conjugacy classes of free factors of rank 2, $OF_2$ \cite{Free_factor_complex_rigidity}, which appears 
in any higher rank complex of conjugacy classes of free factors. This makes studying the theory of the Farey graph essential 
for studying the free factor complex and the complex of conjugacy classes of free factors from the model-theoretic perspective.



The main purpose of this article is to prove the following:

\begin{theorem} \label{main_th}
    The theory of the Farey Graph is axiomatized by
\begin{enumerate}
\item every vertex is contained in an edge;
\item every edge is contained in two triangles; and

\item every finite subgraph of cardinality at least two has two removable vertices\footnote{In fact, it suffices to ask for one removable vertex.}
\end{enumerate}
and   is $\omega$-stable of Morley rank~$\omega$.
\end{theorem}

Note that $\omega$-stability was already  proved  in \cite{modeltheorycurvegraph} for a much wider class of graphs, but without computing the Morley rank.
Our proof proceeds along the lines of the coresponding results for the free pseudospace \cite{TENT_2014} and open generalized polygons \cite{AmmerTent2024}.



\section{The theory of the Farey graph}



The Farey graph is a planar graph that can be defined in a number of different ways. We will use the following
combinatorial definition of the Farey graph, see e.g. \cite{office_hours} and \cite{KURKOFKA2021223}.

\begin{definition}
    The Farey graph $F_1=F_1(e_0)$ of level $1$ consists of two triangles sharing the edge $e_0$.
    Inductively, the Farey graph $F_{n+1}(e_0)$ of level $n + 1$ is the graph obtained from $F_n$ by adding a new vertex $v_e$ for every boundary edge $e$ of $F_{n}$,
    with new edges joining $v_e$ precisely to the vertices of $e$.


    Finally, the \textbf{Farey graph} is the union of these  graphs \[G_F =G_F(e_0):= \bigcup_{n \in \mathbb{N}} F_n.\]

\end{definition}

By construction, every edge in $G_F$ is contained in exactly two triangles and the automorphism group of the Farey graph acts transitively on the set of ordered triangles. 


Note that the symmetry of $F_1(e_0)$ which swaps the two triangles extends to  a symmetry of $F_n, n\in N$  swapping the two hemispheres of $F_n(e_0)$ and fixing $e_0$.
    

We  recall some graph-theoretic notions. 
For $a$ and $b$ in a connected graph $A$, the distance $d(a, b)$ between $a$ and $b$ is the smallest number 
$m$ for which there is a path $a = a_0, a_1, ..., a_m = b$ with $a_i$ in $A$, where 
$a_i$ and $a_{i+1}$ are adjacent for $0 \leq i < m$ and we call a path witnessing this distance a \emph{geodesic}. A vertex $x\in A$ is a cut-vertex in $A$ if $A\setminus\{x\}$ is not connected.
%We may write $d_A(a, b)$ to emphasize the dependence on the graph $A$.

For a vertex $a \in A$ the valency of $a$ in $A$,
 $\val(a)=val_A(a)$, is the number of neighbours of $a$ in $A$.
 
  

It is easy to see from the construction that any finite subgraph $A$ of $G_F$ with $|A|\geq 2$ has (at least) two \emph{removable vertices}
in the sense of the following definition:

\begin{definition}\label{def:removable}
\begin{enumerate}
    \item 
    Let $A$ be a graph, $x\in A$ a vertex. We call $x$ removable from $A$ if $\val(x)$ is at most $1$ or if $\val(x)=2$ and the two neighbours of $x$ form an edge.

    \item If $A$ is a subgraph of $B$, we say that $A$ is strong in $B$ ($A\leq B$) if any non-empty finite subgraph $B_0$ of $B$ not contained in $A$ has a removable vertex $x\in B_0\setminus A$. 
    \end{enumerate}
\end{definition}



\begin{definition}\label{def:K}
    
Let $(\K,\leq)$ be the class of finite graphs $A$ such that
\begin{enumerate}
    \item every edge is contained in at most two triangles, and
    \item every subgraph $A_0$ of $A$ with $|A_0|\geq 2$ has at least two removable vertices.
\end{enumerate}
\end{definition}


\begin{lemma}\label{lem:FG in K}
Every finite subgraph of the Farey graph belongs to $\K$.
\end{lemma}
\begin{proof}
Let $A$ be a finite subgraph of the Farey Graph. Since  condition (1) is clear, it suffices to show that if $|A|\geq 2$, then $A$ has two removable vertice. We may assume that some edge $e$ of $A$ is contained in two triangles (otherwise the claim holds trivially). Suppose $k$ is minimal with $A\subset F_k(e)$. Then in each of the two hemispheres of $F_k(e)$ the last vertex of $A$ added in the construction of $F_k(e)$ is removable.
\end{proof}

\begin{remark} Note that $\emptyset\leq A$ for all $A\in \K$
and that $\leq$ is transitive on $\K$, i.e. if $A\leq B$ and $B\leq C$, then $A\leq C$.
\end{remark} 



The construction of the Farey graph and these definitions immediately imply:

\begin{lemma}\label{lem:strong after removing point}
\begin{enumerate}
\item   If $A$ is strong in $B$, $x,y\in A$, then $d_A(x,y)=d_B(x,y)$.
\item   If $A, B\in \K$, then $A$ is strong in $B$ if and only if we can obtain $B$ from $A$ by successively adjoining removable vertices.
\item  Suppose $A\leq B$ and $x\in B$ is removable over $A$. If one of the following holds
\begin{itemize}
\item $x$ has exactly two neighbours in $A$;
\item $x$ has exactly one neighbour in $A$ and if $val_B(x)\geq 2$, then $x$ is a cut-vertex in $B$; or
\item the connected component of $x$  in $B$ does not intersect $A$.
\end{itemize}
then $A\cup \{x \}\leq B$.
\end{enumerate}
  
  

   
\end{lemma}

\begin{definition}\label{def:free amalgam}
For graphs $A, B, C$ with $A\subset B, C$, we denote by
\[B\otimes_A C\]
the free amalgam of $B$ and $C$ over $A$, i.e. the graph with vertex set $A\cup (B\setminus A)\cup (C\setminus A)$ with the induced edges on $B$ and $C$, respectively, and no edges between vertices of $B\setminus A$ and $C\setminus A$.
\end{definition}

The following is now obvious:
\begin{proposition}\label{prop:amalgamtion}
    The class $(\K,\leq)$ is closed under (AP). More precisely, for  $A, B, C\in \K$ with $A\leq B, C$ there exists $D\in K$ such that $B,C\leq D$ and the embeddings commute.
    
    We let $FG$ denote the Hrushvovki limit of this class.
\end{proposition}
\begin{proof}
Suppose $A, B, C\in \K$ with $A\leq B, C$. We may assume inductively that $|B\setminus A|=1$. Then the free amalgam $D=B\otimes _A C$ belongs to $\K$ unless the vertex $b\in B\setminus A$ has valency $2$ and the edge between the neighbours of $b$ is already contained in two triangles in $C$. In this case, we put $D=C$. Clearly, in either case we have $B, C\leq D$. 
\end{proof}

\begin{remark}
$FG$ and $G_F$ have unbounded diameter.
\end{remark}

We now give axioms for the theory $T=Th(G_F)$ of the Farey graph $G_F$ in the language $L = \{ E \} $ of graphs.

\begin{definition}
    Let $T$ be the theory of graphs axiomatized by 
\begin{enumerate}
   \item every vertex is contained in an edge,
    \item any edge is contained in exactly two triangles; and
    \item every finite subgraph of cardinality at least two has two removable vertices.
\end{enumerate}
\end{definition}
We claim that $T$ is in fact a complete theory.
Clearly, the Farey graph $G_F$ and  the Hrushovski limit $FG$ of $(\K,\leq)$ are models of $T$. 

\begin{remark}\label{rem:triangulated}
Note that (3) implies that any simple cycle in any model of $T$ is triangulated as all vertices in such a cycle have valency $2$.
\end{remark}

We first note the following:

\begin{remark}\label{rem: tree of Fareys}
Let $\G$ be a cycle free graph. Replace every vertex $x\in\G$ by a copy $F_x$ of the Farey graph. For every edge $(x,y)\in \G$, pick vertices $v_x\in F_x, v_y\in F_y$ and identify them. The resulting graph is a model of $T$. 
\end{remark}


We first establish:
\begin{proposition}\label{prop:K-saturated}
    A model  $M$ of $T$ is $\K$-saturated if and only if $\omega$-saturated.
\end{proposition}

\begin{proof}
First suppose $M$ is an $\omega$-saturated model of $T$. We have to show that  $M$ is $\K$-saturated, i.e. we have to show that if $A$ is a finite subgraph of $M$, then $A\in \K$ (which is clear from the axioms of $T$) and that for every pair $A, B\in \K$ with $A\leq B$ and  every strong embedding of $A$ into $M$ we can find a copy of $B$ over $A$ in $M$. (Note that we identify $A$ with its image in $M$.)

Clearly, it suffices to prove this inductively in the case where $B=A\cup\{x\}$ where $x$ is a removable vertex in $B$. 

First consider the case $\val(x)=0$. Note that an $\omega$-saturated model has infinitely many connected components. Since $A$ is finite, we can find a vertex $b\in M$ whose connected component does not intersect $A$. By Lemma~\ref{lem:strong after removing point} we see that if $A\leq M$, then also $A\cup\{ b\}\leq M$.

If $\val(x)=1$, let $a\in A$ be the unique neighbour of $x$. Since every vertex in $M$ has infinite valency and $M$ is $\omega$-saturated, the examples in Remark~\ref{rem: tree of Fareys} show that  we can find a neighbour $b\in M$ of $a$ which is a cut-vertex in $M$.  By Lemma~\ref{lem:strong after removing point} we see that if $A\leq M$, then also $A\cup\{ b\}\leq M$.

    If $x$ has valency $2$ with neighbours $y_1,y_2$ forming an edge in $A$, then this edge is contained in at most $2$ triangles in $B$. Since every edge of $M$ is contained in exactly two triangles, we can find the corresponding triangle in $M$ and we see as before that this extension of $A$ is strong in $M$.

    For the other direction, suppose that $M$ is $\K$-saturated. Then clearly, $M$ is itself a  model of $T$. If $M'$ is an $\omega$-saturated structure elementarily equivalent to $M$, then by the first direction we see that $M'$ is $\K$-saturated. Now $M$ and $M'$ are also partially isomorphic and so $M$ is $\omega$-saturated.
\end{proof}

\begin{corollary}\label{cor:th(FG)}
   $T$ is complete and, hence, is the theory of the Farey Graph.  
\end{corollary}
\begin{proof}
 Clearly, the Farey graph is a model of $T$. To see that $T$ is complete it  is enough to show that  any model of $T$ is elementarily equivalent to $FG$. So let $M\models T$ and let $M'\equiv M$ be $\omega$-saturated. Then by  Proposition~\ref{prop:K-saturated}, $M'$ is $\K$-saturated and so $M'$ and $FG$ are partially isomorphic and hence elementarily equivalent. 
\end{proof}


We can now classify the models of $T$ and we will show that every model is of the form as described in Remark\ref{rem: tree of Fareys} we introduce the following equivalence relation: 

\begin{definition}
    Let $M$ be a  model of $T$.
   Two edges in $M$ are \textbf{equivalent} if they generate the same Farey graph as a subgraph of $M$. 
\end{definition}


\begin{remark}
    Removing a single vertex from $G_F$ leaves the remaining graph connected. Hence for every $x, y \in G_F$, there is a simple cycle in $G_F$ containing $x$ and $y$. Thus, we could equivalently  define the equivalence relation on the edges by saying that two edges are equivalent if and only if they are contained in a simple cycle.
\end{remark}

\begin{lemma}\label{cluster_lemma}
    Two equivalence classes intersect at most in a single vertex.
    
\end{lemma}

\begin{proof}
    If $x$ and $y$ are vertices belonging to equivalence classes $E_1$ and $E_2$, then they
    are connected by paths in $E_1$ and in $E_2$. Either the paths agree, and
    then $E_1$ and $E_2$ contain a common edge and hence $E_1 = E_2$, or we obtain a simple
    cycle. The cycle is triangulated by Remark~\ref{rem:triangulated} and hence again $E_1$ and $E_2$ contain a common edge.
\end{proof}
 
By abuse of terminology we say that a vertex belongs to an equivalence class of edges if it is contained in one of the edges of the equivalence class.

\begin{definition}\label{def:tree}
For any model $M$ of $T$ we define a bipartite graph $G_{tree}(M)$ on the set of vertices of $M$
and the set $\E$ of equivalence classes of $M$, where edges between a vertex $x\in M$ and an equivalence class $E\in\E$ are given by containment,
i.e. a vertex is connected to an equivalence class of edges if it is contained in the equivalence class.
\end{definition}

\begin{lemma}
    For every model $M$ of $T$, $G_{tree}(M)$ is a forest.
\end{lemma}

\begin{proof}
    If the graph $G_{tree}(M)$ contains a cycle, this cycle translates into a cycle in $M$ (by connecting the vertices  of
    the bipartite graph inside their equivalence classes). Any cycle in $M$ is
    triangulated by Remark~\ref{rem:triangulated}, hence the cycle is contained in a single equivalence class.    
\end{proof}

We therefore established:

\begin{theorem}\label{thm:classification}
Every model of $T$ arises from a cycle free graph as described in Remark~\ref{rem: tree of Fareys}.
\end{theorem}

\begin{thebibliography}{999}

 
 \bibitem{AmmerTent2024} A. Ammer, K. Tent, \emph{On the Model Theory of Open Generalized Polygons}. Model Theory, vol. 3, no. 2, pp. 433--448, 2024. 


\bibitem{Free_factor_complex_rigidity} M. Bestvina, M. R Bridson, \emph{Rigidity of the free factor complex}. \textit{arXiv:2306.05941}, 2023. 

\bibitem{office_hours} M. Clay and D. Margalit, \emph{Office Hours with a Geometric Group Theorist}. Princeton University Press, Princeton, NJ, 2017; MR3645425


\bibitem{modeltheorycurvegraph} V. Disarlo, T. Koberda, J. de la Nuez González, \emph{The model theory of the curve graph}. \textit{arXiv:2008.10490}, 2023.

\bibitem{KURKOFKA2021223} 
J. Kurkofka, 
\emph{The Farey Graph is Uniquely Determined by Its Connectivity}. 
Journal of Combinatorial Theory, Series B, 
vol. 151, pp. 223–234, 2021. 

\bibitem{TZ} K. Tent, M. Ziegler, \emph{A course in model theory}. Lecture Notes in Logic, 40. Association for Symbolic Logic, La Jolla, CA; Cambridge University Press, Cambridge, 2012
		
\bibitem{TENT_2014} K. Tent, \emph{The free pseudospace is n-ample, but not (n + 1)-ample}. The Journal of Symbolic Logic, 
vol. 79, no. 2, pp. 410–428, 2014. 

\end{thebibliography}
\end{document}
